\section{What the public timetable does and does not establish}\label{sec:public}
The analytical construction proves possibility, not operational frequency.
The public case uses the Hildenbrand timetable from the Sistig data
\citep{sistig2025dataset,sistig2025electrification}: 37 mandatory services
and the complete directed deadhead matrix. Depot 15 and depot 16 define
separate one-depot modifications of the original network. Source timetable
offsets, including next-day services, are retained; source data do not
provide the synthetic energy and cost assumptions below.

Service consumption is 1.20 kWh/km plus 16 kWh per service hour; directed
deadheads use 0.96 kWh/km plus 16 kWh per travel hour, and explicit off-depot
waiting uses 16 kWh/h. Depot-parked auxiliary draw is zero. A used bus has
400 kWh usable inventory, zero reserve, unit efficiency and 360 kW acceptance.
One 360 kW connector is shared. Every used bus starts full and returns full
by a synthetic 30:00 service-block boundary. There is no taper or switching
time. Monetary movement cost is zero; movement time and energy remain
physical constraints. The used-bus fee is 100. This finite block is not a
certificate of repeating daily operation.

The initial flat tariff is 0.2 per kWh. Because the complete-plan objective
is linear, its physical and hull optima have equal value regardless of the
discrete assignments. A native numerical flat optimum near 408.53 at depot 15
is tolerance-qualified and does not identify the ideal stored-input rational
optimum. Exact rational evidence instead provides a two-bus physical witness
and a cardinality-gated path-cover relaxation, giving outward bounds
\[
 D^{\rm flat}_{15}\in[404.92,408.54].
\]
An exact one-bus obstruction establishes minimum fleet size two. Depot 16
has an exact lower bound 414.39, but its reported two-bus upper witness is
numerical; the resulting interval must be labeled mixed evidence.

On the same finite depot 15 ideal model, let
\begin{equation}\label{eq:publicF}
 F(e)=\sum_{t=1}^{30}\left(0.2e_t+\frac{e_t^2}{1800}\right).
\end{equation}
The displayed coefficients denote the original stored binary values in the
exact certificates.\footnote{Precisely, the linear coefficient is
$3602879701896397/18014398509481984$ and the quadratic Hessian coefficient is
$5124095576030431/4611686018427387904$; the quadratic term is half the latter
times $e_t^2$. The familiar decimal expression displays their intended scale.}
A new exact two-bus witness and an independently reconstructed rational
lower certificate for a per-leg perspective energy-flow relaxation
give the presently established outward bounds
\begin{equation}\label{eq:publicbounds}
 434.17\le\CH_{15}\le D_{15}\le441.54,\qquad
 0\le D_{15}-\CH_{15}\le7.37.
\end{equation}
The new exact lower is approximately $434.17948$, rounded downward here;
its validity follows from a supply tangent, sign-valid rational row
multipliers and finite-bound compensation for stationarity residuals.
It improves the earlier energy-floor lower of 429.22. The continuous
energy-flow relaxation is not asserted to equal the complete-fleet hull.
A separate native integer solve returned OPTIMAL with objective
$441.532912\ldots$, but its candidate failed strict physical replay and
did not itself improve the exact upper or establish ideal-model optimality.
The native status and the independently checked rational lower are
different evidence. The new admitted witness has cost
$442.123049\ldots$. A time-limited classical MIQP supplied its routes and
charging proposal; two terminal energy changes smaller than $7\cdot10^{-12}$
kWh restored exact full inventories. A separate rational event replay then
verified all services, directed movements, battery bounds, shared connector
use and full replenishment. The original unadmitted numerical proposal is
retained separately. No route or nonterminal charging allocation changed.

For any fleet with at least three buses, its operating cost and only the
linear mandatory-service-energy component already exceed $477.82$. Movement
consumption and the quadratic cost are nonnegative. Such a fleet cannot
improve the new physical upper. Together with the one-bus obstruction, this
establishes exactly two buses at any physical planner optimum. It does not
exclude larger-fleet components from arbitrary hull mixtures or from a
price-response problem at another tariff.

Separately, a second strictly admitted physical plan and the protected witness
give a feasible complete-plan mixture with value $441.600681\ldots$, hence
$\CH_{15}\le441.61$. Its negative initial directional derivative proves that
the protected witness is unsupported at its own marginal prices. A mixture
is not a dispatch, and an unsupported incumbent is not proof that an optimal
physical fleet is unsupported. A subsequent fixed-route charging QP returned
the same lower numerical value $441.532912\ldots$, but its first exact dual
recovery failed on a multiplier sign; the conditional global own-price
checks were therefore unattempted. A distinct rational phase-I method then
constructed a new exact primal, admitted it by independent physical replay,
and found sign-valid multipliers certifying its fixed-route charging optimum.
This supplies the stronger physical upper in \eqref{eq:publicbounds}; it also
supersedes the earlier mixture upper as a bound on $\CH$. The original failure
remains recorded. Fixed-route optimality still does not exclude a profitable
complete-fleet deviation or certify global marginal-price support.
The bounds neither determine
$D_{15}$ or $\CH_{15}$ nor establish a positive gap. No new public sign
certificate is asserted in this draft.

There are distinct valid ways to close this question. An exact physical
witness $x$ and a globally valid linear-response lower bound matching
$c(x)+\nabla F(e(x))^\top e(x)$ prove support and hence $D=\CH$.
Alternatively, a physical lower bound strictly above a feasible mixture's
hull upper bound proves positive gap. A globally certified hull optimizer
can also prove failure through nonimplementability: strict convexity in
\eqref{eq:publicF} fixes its load and operating cost, and a hull operating
cost not in $100\mathbb Z_+$ cannot be the cost of a physical fleet.
All these routes require a global certificate or a complete infeasibility
argument. A finite plan bank, incumbent regret, or a small native solver
gap does not settle the sign of the ideal-model gap.
